\section{Discussion}
\label{sec:discussion}

The matched comparisons yield three practical findings: the prediction rule
produces the largest observed component difference, KD primarily affects
retention, and total exemplar memory has a larger observed effect than the
tested retrieval range.

\subsection{What Drives Performance}
\label{subsec:what-drives-performance}

We find that the final prediction rule is the largest measured lever in this
protocol. Replacing nearest-mean exemplar (NME) prediction with head-logit
prediction reduces final average accuracy by $10.38$ pp and increases
forgetting by $32.50$ pp (Table~\ref{tab:t2-component-ablations}). Both variants
use the same training procedure and differ only at prediction time; under this
comparison, head-logit prediction underperforms the selected-exemplar prototypes
used by NME.

KD has a distinct role. Removing it changes average accuracy by $-0.21$ pp
but raises forgetting by $10.93$ pp, consistent with KD acting as a retention
mechanism. Herding also matters once budget and replay procedure are held
fixed: replacing it with random selection reduces average accuracy by $4.83$
pp and increases forgetting by $4.50$ pp. By contrast, changing the training
head from cosine-margin to linear under NME prediction produces smaller
observed changes ($-1.42$ pp accuracy and $+0.83$ pp forgetting). Within this
protocol, prediction and exemplar selection therefore deserve priority over
head geometry when the objective is to improve the accuracy--forgetting trade-off.

\subsection{Resource Allocation and Design Implications}
\label{subsec:resource-allocation}

The resource sweep shows that storage matters more than retrieval count in the
tested range. Reducing memory from 2,000 to 500 exemplars loses $8.17$ pp in
accuracy and adds $5.77$ pp of forgetting. Doubling memory from 2,000 to
4,000 improves both metrics, whereas doubling retrieval from 64 to 128 does
not improve either mean metric (Table~\ref{tab:t3-resource-sensitivity}).
Retrieval may still matter, but expanding storage is the higher-priority
allocation in this setting.

For the evaluated setting, a replay system should allocate memory across all
observed classes, rebuild exemplars with herding, train with KD, and use NME
prediction. When additional resources are available, the measurements favor
expanding stored memory before increasing retrieval from 64 samples.

\subsection{Limitations}
\label{subsec:limitations}

Our claims are bounded by the protocol's coverage. All experiments use
CIFAR-100, one 10-task class partition induced by split seed 13, and a
ResNet-18 backbone. The three seeds quantify training variation but do not
measure sensitivity to alternative class orders, datasets, architectures, or
task granularities. The paired comparisons are based on three matched seeds;
they support the reported directions and magnitudes but do not provide
high-power estimates of small effects.

The memory sweep covers 500, 2,000, and 4,000 exemplars. Although the increase
from 2,000 to 4,000 is smaller than the loss from 2,000 to 500, these three
points do not establish a saturation threshold or a scaling law. The baseline
comparison isolates replay, selection, and prediction choices under a common
protocol. Extending the same matched ablations across class orders and standard
benchmarks is necessary to test whether the observed component ranking transfers.
